\documentclass[11pt,a4paper]{article}
\usepackage[utf8]{inputenc}
\usepackage[T1]{fontenc}
\usepackage[brazil]{babel}
\usepackage{lmodern}
\usepackage{geometry}
\geometry{left=2.2cm,right=2.2cm,top=2.1cm,bottom=2.2cm}
\usepackage{amsmath,amssymb}
\usepackage{booktabs,array,longtable,calc}
\usepackage{xcolor}
\usepackage{tikz}
\usetikzlibrary{arrows.meta,positioning}
\usepackage{fancyhdr}
\usepackage{hyperref}
\definecolor{accent}{HTML}{203D58}
\definecolor{soft}{HTML}{EDF2F6}
\hypersetup{colorlinks=true,linkcolor=accent,urlcolor=accent,
  pdftitle={Inteligência Artificial - Resolução da Questão 1},
  pdfsubject={Rede bayesiana para diagnóstico de um farol de bicicleta}}
\pagestyle{fancy}
\fancyhf{}
\fancyhead[L]{\small\textcolor{accent}{Inteligência Artificial}}
\fancyhead[R]{\small Questão 1}
\fancyfoot[C]{\small\thepage}
\setlength{\headheight}{14pt}
\setlength{\parindent}{0pt}
\setlength{\parskip}{5pt}
\setlength{\tabcolsep}{8pt}
\renewcommand{\arraystretch}{1.12}
\newcommand{\snow}{\texttt{snow\_covered}}
\newcommand{\itemtitle}[2]{\section*{\textcolor{accent}{Item (#1) --- #2}}}
\newcommand{\subtitulo}[1]{\par\medskip\textbf{\textcolor{accent}{#1}}\par}

\setcounter{secnumdepth}{0}

\begin{document}
{\LARGE\bfseries\textcolor{accent}{Raciocínio probabilístico}}\\[4pt]
{\large Resolução da questão 1: diagnóstico de um farol a dínamo}

\textbf{Disciplina:} Inteligência Artificial \hfill \textbf{2\textsuperscript{o} Trabalho}\par
\textbf{Professor:} Edjard Mota \hfill \textbf{Data do enunciado:} 11/10/2026

\itemtitle{a}{Rede bayesiana causal}\label{item-a}

\subtitulo{Variáveis}\label{variuxe1veis}

{\def\LTcaptype{none} % do not increment counter
\begin{longtable}[]{@{}lll@{}}
\toprule\noalign{}
Variável & Significado & Valores \\
\midrule\noalign{}
\endhead
\bottomrule\noalign{}
\endlastfoot
\(Str\) & Condição da rua & \texttt{dry}, \texttt{wet}, \texttt{snow\_covered} \\
\(Flw\) & Volante do dínamo desgastado & \(t\), \(f\) \\
\(R\) & Dínamo deslizando & \(t\), \(f\) \\
\(V\) & Dínamo apresenta tensão & \(t\), \(f\) \\
\(B\) & Lâmpada em boas condições & \(t\), \(f\) \\
\(K\) & Cabo em boas condições & \(t\), \(f\) \\
\(Li\) & Luz ligada & \(t\), \(f\) \\
\end{longtable}
}

Em todas as tabelas, \(t\) significa verdadeiro e \(f\), falso.

\subtitulo{Construção da rede}\label{construuxe7uxe3o-da-rede}

A condição da rua e o desgaste do volante afetam a aderência do dínamo. O deslizamento afeta a geração de tensão. A luz depende da tensão, da lâmpada e do cabo. Assim, adotamos a seguinte rede:

\begin{center}
\begin{tikzpicture}[
causal/.style={-{Stealth[length=2.5mm]},thick,draw=accent},
var/.style={draw=accent,rounded corners=3pt,fill=soft,minimum width=23mm,minimum height=9mm},
every node/.style={align=center}]
\node[var] (str) at (-2.2,3) {$Str$\\{\scriptsize condição da rua}};
\node[var] (flw) at (2.2,3) {$Flw$\\{\scriptsize volante desgastado}};
\node[var] (r) at (0,1.85) {$R$\\{\scriptsize dínamo deslizando}};
\node[var] (v) at (0,0.65) {$V$\\{\scriptsize presença de tensão}};
\node[var] (b) at (-3.2,0.65) {$B$\\{\scriptsize lâmpada ok}};
\node[var] (k) at (3.2,0.65) {$K$\\{\scriptsize cabo ok}};
\node[var] (li) at (0,-0.65) {$Li$\\{\scriptsize luz ligada}};
\draw[causal] (str)--(r);
\draw[causal] (flw)--(r);
\draw[causal] (r)--(v);
\draw[causal] (v)--(li);
\draw[causal] (b)--(li);
\draw[causal] (k)--(li);
\end{tikzpicture}
\end{center}

\clearpage
O conjunto de arestas é

\[
E=\{Str\to R,\ Flw\to R,\ R\to V,\ V\to Li,\ B\to Li,\ K\to Li\}.
\]

{\def\LTcaptype{none} % do not increment counter
\begin{longtable}[]{@{}ll@{}}
\toprule\noalign{}
Nó & Pais \\
\midrule\noalign{}
\endhead
\bottomrule\noalign{}
\endlastfoot
\(Str\), \(Flw\), \(B\), \(K\) & Nenhum \\
\(R\) & \(Str\), \(Flw\) \\
\(V\) & \(R\) \\
\(Li\) & \(V\), \(B\), \(K\) \\
\end{longtable}
}

As condições \(P(V\mid R,Str)=P(V\mid R)\) e \(P(V\mid R,Flw)=P(V\mid R)\) são compatíveis com \(R\) como único pai de \(V\). Conhecido o deslizamento, rua e desgaste não acrescentam informação sobre a tensão no modelo. A condição \(P(Li\mid V,R)=P(Li\mid V)\) é compatível com o efeito de \(R\) sobre a luz sendo intermediado por \(V\). A tabela de \(Li\) fornecida no enunciado utiliza justamente \(V\), \(B\) e \(K\) como pais.

\subtitulo{Hipótese de modelagem e fatoração}\label{hipuxf3tese-de-modelagem-e-fatorauxe7uxe3o}

Adotei a interpretação causal descrita e a hipótese usual de \textbf{independência conjunta} dos nós de origem \(Str\), \(Flw\), \(B\) e \(K\). Essa hipótese é compatível com a independência aos pares do enunciado, mas é mais forte: independência aos pares, sozinha, não implica independência conjunta nem determina uma única rede.

A distribuição conjunta da rede escolhida se decompõe em

\[
\begin{aligned}
P(Str,Flw,R,V,B,K,Li)
={}&P(Str)P(Flw)P(B)P(K)\\
&\cdot P(R\mid Str,Flw)P(V\mid R)P(Li\mid V,B,K).
\end{aligned}
\]

Sem observar a luz, os caminhos que ligam \(B\) ou \(K\) à cadeia \(Str,Flw,R,V\) passam pelo colisor \(Li\) e ficam bloqueados. Isso explica as independências marginais \((R,B)\), \((R,K)\), \((V,B)\) e \((V,K)\).

\textbf{Resposta:} a rede tem sete nós e as seis arestas indicadas no diagrama; não há ciclos dirigidos.

\clearpage
\itemtitle{b}{Tabelas de probabilidades}\label{item-b}

Cada nó precisa de uma distribuição condicionada aos seus pais. Para um nó sem pais, essa tabela é sua distribuição a priori.

\begin{center}
\begin{tikzpicture}[
causal/.style={-{Stealth[length=2.5mm]},thick,draw=accent},
var/.style={draw=accent,rounded corners=3pt,fill=soft,minimum width=23mm,minimum height=9mm},
every node/.style={align=center}]
\node[var] (str) at (-2.2,3) {$Str$\\{\scriptsize $P(Str)$}};
\node[var] (flw) at (2.2,3) {$Flw$\\{\scriptsize $P(Flw)$}};
\node[var] (r) at (0,1.85) {$R$\\{\scriptsize $P(R\mid Str,Flw)$}};
\node[var] (v) at (0,0.65) {$V$\\{\scriptsize $P(V\mid R)$}};
\node[var] (b) at (-3.2,0.65) {$B$\\{\scriptsize $P(B)$}};
\node[var] (k) at (3.2,0.65) {$K$\\{\scriptsize $P(K)$}};
\node[var] (li) at (0,-0.65) {$Li$\\{\scriptsize $P(Li\mid V,B,K)$}};
\draw[causal] (str)--(r);
\draw[causal] (flw)--(r);
\draw[causal] (r)--(v);
\draw[causal] (v)--(li);
\draw[causal] (b)--(li);
\draw[causal] (k)--(li);
\end{tikzpicture}
\end{center}

Há \textbf{sete tabelas}. A CPT de \(Li\) já consta do enunciado; faltam as seis distribuições de \(Str\), \(Flw\), \(B\), \(K\), \(R\) e \(V\). Abaixo estão suas formas gerais. As escolhas numéricas são desenvolvidas no \hyperref[item-c]{item (c)}.

\subtitulo{\texorpdfstring{1. Distribuição de \(Str\)}{1. Distribuição de Str}}\label{distribuiuxe7uxe3o-de-str}

{\def\LTcaptype{none} % do not increment counter
\begin{longtable}[]{@{}lr@{}}
\toprule\noalign{}
\(Str\) & Probabilidade \\
\midrule\noalign{}
\endhead
\bottomrule\noalign{}
\endlastfoot
\texttt{dry} & \(p_d\) \\
\texttt{wet} & \(p_w\) \\
\texttt{snow\_covered} & \(p_s\) \\
\end{longtable}
}

Os parâmetros devem satisfazer \(p_d,p_w,p_s\geq 0\) e \(p_d+p_w+p_s=1\).

\subtitulo{\texorpdfstring{2--4. Distribuições de \(Flw\), \(B\) e \(K\)}{2--4. Distribuições de Flw, B e K}}\label{distribuiuxe7uxf5es-de-flw-b-e-k}

Definimos \(\alpha=P(Flw=t)\), \(\beta=P(B=t)\) e \(\gamma=P(K=t)\).

{\def\LTcaptype{none} % do not increment counter
\begin{longtable}[]{@{}lrr@{}}
\toprule\noalign{}
Nó & \(P(t)\) & \(P(f)\) \\
\midrule\noalign{}
\endhead
\bottomrule\noalign{}
\endlastfoot
\(Flw\) & \(\alpha\) & \(1-\alpha\) \\
\(B\) & \(\beta\) & \(1-\beta\) \\
\(K\) & \(\gamma\) & \(1-\gamma\) \\
\end{longtable}
}

Cada parâmetro pertence ao intervalo \([0,1]\).

\clearpage
\subtitulo{\texorpdfstring{5. CPT de \(R\)}{5. CPT de R}}\label{cpt-de-r}

Definimos \(r_{s,w}=P(R=t\mid Str=s,Flw=w)\).

{\def\LTcaptype{none} % do not increment counter
\begin{longtable}[]{@{}llrr@{}}
\toprule\noalign{}
\(Str\) & \(Flw\) & \(P(R=t\mid Str,Flw)\) & \(P(R=f\mid Str,Flw)\) \\
\midrule\noalign{}
\endhead
\bottomrule\noalign{}
\endlastfoot
\texttt{dry} & \(f\) & \(r_{dry,f}\) & \(1-r_{dry,f}\) \\
\texttt{dry} & \(t\) & \(r_{dry,t}\) & \(1-r_{dry,t}\) \\
\texttt{wet} & \(f\) & \(r_{wet,f}\) & \(1-r_{wet,f}\) \\
\texttt{wet} & \(t\) & \(r_{wet,t}\) & \(1-r_{wet,t}\) \\
\texttt{snow\_covered} & \(f\) & \(r_{snow,f}\) & \(1-r_{snow,f}\) \\
\texttt{snow\_covered} & \(t\) & \(r_{snow,t}\) & \(1-r_{snow,t}\) \\
\end{longtable}
}

O índice \(snow\) abrevia o valor \texttt{snow\_covered}. São \(3\times2=6\) combinações dos pais e seis parâmetros livres, todos em \([0,1]\).

\subtitulo{\texorpdfstring{6. CPT de \(V\)}{6. CPT de V}}\label{cpt-de-v}

Definimos \(v_r=P(V=t\mid R=r)\).

{\def\LTcaptype{none} % do not increment counter
\begin{longtable}[]{@{}lrr@{}}
\toprule\noalign{}
\(R\) & \(P(V=t\mid R)\) & \(P(V=f\mid R)\) \\
\midrule\noalign{}
\endhead
\bottomrule\noalign{}
\endlastfoot
\(f\) & \(v_f\) & \(1-v_f\) \\
\(t\) & \(v_t\) & \(1-v_t\) \\
\end{longtable}
}

São duas combinações do único pai e dois parâmetros livres em \([0,1]\).

\subtitulo{\texorpdfstring{7. CPT de \(Li\) fornecida no enunciado}{7. CPT de Li fornecida no enunciado}}\label{cpt-de-li-fornecida-no-enunciado}

Interpretamos a coluna \(P(Li)\) da tabela original como \(P(Li=t\mid V,B,K)\). Preservamos os oito valores fornecidos e acrescentamos os complementos para \(Li=f\).

{\def\LTcaptype{none} % do not increment counter
\begin{longtable}[]{@{}lllrr@{}}
\toprule\noalign{}
\(V\) & \(B\) & \(K\) & \(P(Li=t\mid V,B,K)\) & \(P(Li=f\mid V,B,K)\) \\
\midrule\noalign{}
\endhead
\bottomrule\noalign{}
\endlastfoot
\(t\) & \(t\) & \(t\) & \(0{,}99\) & \(0{,}01\) \\
\(t\) & \(t\) & \(f\) & \(0{,}01\) & \(0{,}99\) \\
\(t\) & \(f\) & \(t\) & \(0{,}01\) & \(0{,}99\) \\
\(t\) & \(f\) & \(f\) & \(0{,}001\) & \(0{,}999\) \\
\(f\) & \(t\) & \(t\) & \(0{,}30\) & \(0{,}70\) \\
\(f\) & \(t\) & \(f\) & \(0{,}005\) & \(0{,}995\) \\
\(f\) & \(f\) & \(t\) & \(0{,}005\) & \(0{,}995\) \\
\(f\) & \(f\) & \(f\) & \(0\) & \(1\) \\
\end{longtable}
}

São \(2^3=8\) combinações dos pais. Para qualquer variável booleana \(X\), o complemento decorre de

\[
P(X=f\mid\text{pais})=1-P(X=t\mid\text{pais}).
\]

\textbf{Resposta:} as sete tabelas ficam associadas aos nós como no diagrama; suas linhas somam 1. O total de parâmetros livres da estrutura é \(2+1+1+1+6+2+8=21\), dos quais oito já foram fixados pela CPT de \(Li\) do enunciado.

\clearpage
\itemtitle{c}{Valores plausíveis}\label{item-c}

A rede adotada é \(Str\to R\leftarrow Flw\), \(R\to V\), \(V\to Li\), \(B\to Li\) e \(K\to Li\), com os nós de origem conjuntamente independentes. A CPT de \(Li\) é preservada.

\subtitulo{Distribuições a priori}\label{distribuiuxe7uxf5es-a-priori}

{\def\LTcaptype{none} % do not increment counter
\begin{longtable}[]{@{}lr@{}}
\toprule\noalign{}
\(Str\) & \(P(Str)\) \\
\midrule\noalign{}
\endhead
\bottomrule\noalign{}
\endlastfoot
\texttt{dry} & \(0{,}60\) \\
\texttt{wet} & \(0{,}30\) \\
\texttt{snow\_covered} & \(0{,}10\) \\
\end{longtable}
}

\[
0{,}60+0{,}30+0{,}10=1.
\]

Escolhi um cenário em que ruas secas são mais frequentes que ruas molhadas e ruas com neve.

{\def\LTcaptype{none} % do not increment counter
\begin{longtable}[]{@{}lrr@{}}
\toprule\noalign{}
Nó & \(P(t)\) & \(P(f)\) \\
\midrule\noalign{}
\endhead
\bottomrule\noalign{}
\endlastfoot
\(Flw\) & \(0{,}20\) & \(0{,}80\) \\
\(B\) & \(0{,}95\) & \(0{,}05\) \\
\(K\) & \(0{,}98\) & \(0{,}02\) \\
\end{longtable}
}

Assim, há 20\% de probabilidade de desgaste do volante, 95\% de a lâmpada estar em boas condições e 98\% de o cabo estar em boas condições.

\subtitulo{\texorpdfstring{CPT de \(R\)}{CPT de R}}\label{cpt-de-r-1}

{\def\LTcaptype{none} % do not increment counter
\begin{longtable}[]{@{}llrr@{}}
\toprule\noalign{}
\(Str\) & \(Flw\) & \(P(R=t\mid Str,Flw)\) & \(P(R=f\mid Str,Flw)\) \\
\midrule\noalign{}
\endhead
\bottomrule\noalign{}
\endlastfoot
\texttt{dry} & \(f\) & \(0{,}02\) & \(0{,}98\) \\
\texttt{dry} & \(t\) & \(0{,}10\) & \(0{,}90\) \\
\texttt{wet} & \(f\) & \(0{,}20\) & \(0{,}80\) \\
\texttt{wet} & \(t\) & \(0{,}50\) & \(0{,}50\) \\
\texttt{snow\_covered} & \(f\) & \(0{,}60\) & \(0{,}40\) \\
\texttt{snow\_covered} & \(t\) & \(0{,}90\) & \(0{,}10\) \\
\end{longtable}
}

Para um volante sem desgaste, o risco de deslizamento aumenta de \(0{,}02\) para \(0{,}20\) e depois para \(0{,}60\) ao passar de rua seca para molhada e com neve. Para um volante desgastado, aumenta de \(0{,}10\) para \(0{,}50\) e depois para \(0{,}90\). Em cada condição de rua, o desgaste aumenta a probabilidade de deslizamento.

\subtitulo{\texorpdfstring{CPT de \(V\)}{CPT de V}}\label{cpt-de-v-1}

{\def\LTcaptype{none} % do not increment counter
\begin{longtable}[]{@{}lrr@{}}
\toprule\noalign{}
\(R\) & \(P(V=t\mid R)\) & \(P(V=f\mid R)\) \\
\midrule\noalign{}
\endhead
\bottomrule\noalign{}
\endlastfoot
\(f\) & \(0{,}98\) & \(0{,}02\) \\
\(t\) & \(0{,}10\) & \(0{,}90\) \\
\end{longtable}
}

Sem deslizamento, o dínamo apresenta tensão com probabilidade de 98\%. Com deslizamento, essa probabilidade cai para 10\%. Não a fixei em zero porque pode haver contato residual ou deslizamento parcial.

\clearpage
\subtitulo{\texorpdfstring{CPT de \(Li\): valores do enunciado}{CPT de Li: valores do enunciado}}\label{cpt-de-li-valores-do-enunciado}

{\def\LTcaptype{none} % do not increment counter
\begin{longtable}[]{@{}lllrr@{}}
\toprule\noalign{}
\(V\) & \(B\) & \(K\) & \(P(Li=t\mid V,B,K)\) & \(P(Li=f\mid V,B,K)\) \\
\midrule\noalign{}
\endhead
\bottomrule\noalign{}
\endlastfoot
\(t\) & \(t\) & \(t\) & \(0{,}99\) & \(0{,}01\) \\
\(t\) & \(t\) & \(f\) & \(0{,}01\) & \(0{,}99\) \\
\(t\) & \(f\) & \(t\) & \(0{,}01\) & \(0{,}99\) \\
\(t\) & \(f\) & \(f\) & \(0{,}001\) & \(0{,}999\) \\
\(f\) & \(t\) & \(t\) & \(0{,}30\) & \(0{,}70\) \\
\(f\) & \(t\) & \(f\) & \(0{,}005\) & \(0{,}995\) \\
\(f\) & \(f\) & \(t\) & \(0{,}005\) & \(0{,}995\) \\
\(f\) & \(f\) & \(f\) & \(0\) & \(1\) \\
\end{longtable}
}

Os valores não nulos com \(V=f\) pertencem à tabela fornecida e não foram substituídos por um modelo determinístico de circuito.

\subtitulo{Normalização e coerência}\label{normalizauxe7uxe3o-e-coeruxeancia}

Todas as probabilidades estão no intervalo \([0,1]\). Em cada linha booleana, a probabilidade de \(f\) é o complemento da probabilidade de \(t\); por exemplo,

\[
P(R=f\mid Str=\text{snow\_covered},Flw=t)=1-0{,}90=0{,}10,
\]

\[
P(Li=f\mid V=t,B=f,K=f)=1-0{,}001=0{,}999.
\]

A independência conjunta das origens garante que as probabilidades de \(B\) e \(K\) não dependam de \(R\) ou de \(V\). As CPTs de \(V\) e de \(Li\) também respeitam as independências condicionais indicadas: \(V\) depende apenas de \(R\), e o efeito de \(R\) sobre \(Li\) passa por \(V\).

\textbf{Resposta:} as seis distribuições faltantes ficam numericamente especificadas, e a sétima permanece igual à do enunciado. São esses valores que serão usados no \hyperref[item-e]{item (e)}.

\clearpage
\itemtitle{d}{\texorpdfstring{ausência de aresta entre \texttt{Str} e \texttt{Li}}{ausência de aresta entre Str e Li}}\label{item-d}

Na rede adotada, a condição da rua altera a probabilidade de deslizamento, que altera a geração de tensão, que influencia a luz:

\[
Str\longrightarrow R\longrightarrow V\longrightarrow Li.
\]

Portanto, há um efeito \textbf{indireto}. Os pais de \(Li\) são exatamente \(V\), \(B\) e \(K\), como também indica a CPT fornecida no enunciado. A rua não é um pai da luz.

\subtitulo{Justificativa por d-separação}\label{justificativa-por-d-separauxe7uxe3o}

O único caminho entre \(Str\) e \(Li\) passa por \(R\) e \(V\). Ao condicionar em \(V\), um nó não colisor nesse caminho, bloqueamos o caminho. Pela propriedade de d-separação da rede,

\[
Str\perp Li\mid V.
\]

Mesmo quando a lâmpada e o cabo também são conhecidos,

\[
P(Li\mid Str,V,B,K)=P(Li\mid V,B,K).
\]

Conhecidos os pais da luz, informar a condição da rua não muda sua probabilidade. A condição \(P(Li\mid V,R)=P(Li\mid V)\) do enunciado também é compatível com essa mediação: ao conhecer a tensão, saber se o dínamo desliza não acrescenta informação sobre a luz.

\subtitulo{Verificação pela fatoração}\label{verificauxe7uxe3o-pela-fatorauxe7uxe3o}

Escreva \(s\), \(v\), \(b\) e \(k\) para valores fixados de \(Str\), \(V\), \(B\) e \(K\). Na fatoração da rede, o único fator que envolve \(Li\) é \(P(Li\mid v,b,k)\). Ao somar os estados não observados de \(Flw\) e \(R\) e normalizar, esse fator pode sair da soma, resultando em

\[
P(Li\mid s,v,b,k)=P(Li\mid v,b,k).
\]

Se \(B\) e \(K\) não forem observados, a independência conjunta das origens permite marginalizar:

\[
P(Li\mid V=v,Str=s)
=\sum_{b\in\{t,f\}}\sum_{k\in\{t,f\}}
P(Li\mid v,b,k)P(B=b)P(K=k).
\]

O lado direito não depende de \(s\) e é igual a \(P(Li\mid V=v)\).

\subtitulo{Conclusão}\label{conclusuxe3o}

Na estrutura causal escolhida, toda influência da rua sobre a luz é transmitida por \(R\) e \(V\). Assim, \textbf{não há aresta direta \(Str\to Li\)}. Também não há aresta inversa \(Li\to Str\), pois o estado da luz não causa a condição da rua.

A ausência da aresta não significa que \(Str\) e \(Li\) sejam independentes sem condicionamento: o caminho indireto continua ativo quando \(V\) não é observado. A conclusão refere-se à rede causal adotada; as independências fornecidas, isoladamente, não demonstram que ela seja a única estrutura possível.

\clearpage
\itemtitle{e}{Tensão com a rua coberta de neve}\label{item-e}

Como \(V\) é booleana, vamos calcular a probabilidade de \(V=t\) e sua complementar para \(V=f\). Abreviamos a evidência \(Str=\text{snow\_covered}\) por \(S\).

Usando as probabilidades escolhidas no \hyperref[item-c]{item (c)}:

\[
P(Flw=t)=0{,}20,\qquad P(Flw=f)=0{,}80,
\]

\[
P(R=t\mid S,Flw=t)=0{,}90,\qquad
P(R=t\mid S,Flw=f)=0{,}60,
\]

\[
P(V=t\mid R=t)=0{,}10,\qquad
P(V=t\mid R=f)=0{,}98.
\]

Pela independência entre \(Str\) e \(Flw\), observar neve não altera a probabilidade de desgaste:

\[
P(Flw=t\mid S)=0{,}20,\qquad P(Flw=f\mid S)=0{,}80.
\]

\subtitulo{1. Marginalizar o desgaste para obter o deslizamento}\label{marginalizar-o-desgaste-para-obter-o-deslizamento}

Pela lei da probabilidade total, somamos os dois estados de \(Flw\):

\[
\begin{aligned}
P(R=t\mid S)
={}&P(R=t\mid S,Flw=t)P(Flw=t\mid S)\\
&+P(R=t\mid S,Flw=f)P(Flw=f\mid S)\\
={}&0{,}90\cdot0{,}20+0{,}60\cdot0{,}80\\
={}&0{,}18+0{,}48\\
={}&0{,}66.
\end{aligned}
\]

Consequentemente,

\[
P(R=f\mid S)=1-0{,}66=0{,}34.
\]

\subtitulo{2. Marginalizar o deslizamento para obter a tensão}\label{marginalizar-o-deslizamento-para-obter-a-tensuxe3o}

Novamente pela lei da probabilidade total,

\[
P(V=t\mid S)=\sum_{r\in\{t,f\}}P(V=t\mid R=r,S)P(R=r\mid S).
\]

O enunciado fornece \(P(V\mid R,Str)=P(V\mid R)\). Portanto, conhecido \(R\), podemos retirar a condição sobre a rua do primeiro fator:

\[
\begin{aligned}
P(V=t\mid S)
={}&P(V=t\mid R=t)P(R=t\mid S)\\
&+P(V=t\mid R=f)P(R=f\mid S)\\
={}&0{,}10\cdot0{,}66+0{,}98\cdot0{,}34\\
={}&0{,}066+0{,}3332\\
={}&\boxed{0{,}3992=39{,}92\%}.
\end{aligned}
\]

\clearpage
\subtitulo{3. Conferência pelos quatro caminhos}\label{conferuxeancia-pelos-quatro-caminhos}

A mesma soma pode ser aberta nas quatro combinações de desgaste e deslizamento:

{\def\LTcaptype{none} % do not increment counter
\begin{longtable}[]{@{}llr@{}}
\toprule\noalign{}
\(Flw\) & \(R\) & \(P(Flw\mid S)P(R\mid S,Flw)P(V=t\mid R)\) \\
\midrule\noalign{}
\endhead
\bottomrule\noalign{}
\endlastfoot
\(t\) & \(t\) & \(0{,}20\cdot0{,}90\cdot0{,}10=0{,}018\) \\
\(t\) & \(f\) & \(0{,}20\cdot0{,}10\cdot0{,}98=0{,}0196\) \\
\(f\) & \(t\) & \(0{,}80\cdot0{,}60\cdot0{,}10=0{,}048\) \\
\(f\) & \(f\) & \(0{,}80\cdot0{,}40\cdot0{,}98=0{,}3136\) \\
\end{longtable}
}

Logo,

\[
0{,}018+0{,}0196+0{,}048+0{,}3136=0{,}3992.
\]

Essa abertura usa a CPT \(P(V\mid R)\) da rede adotada.

\subtitulo{4. Distribuição final}\label{distribuiuxe7uxe3o-final}

\[
P(V=f\mid S)=1-0{,}3992=\boxed{0{,}6008=60{,}08\%}.
\]

{\def\LTcaptype{none} % do not increment counter
\begin{longtable}[]{@{}lr@{}}
\toprule\noalign{}
Estado de \(V\) & \(P(V\mid Str=\text{snow\_covered})\) \\
\midrule\noalign{}
\endhead
\bottomrule\noalign{}
\endlastfoot
\(t\) --- apresenta tensão & \(0{,}3992\) (\(39{,}92\%\)) \\
\(f\) --- não apresenta tensão & \(0{,}6008\) (\(60{,}08\%\)) \\
\end{longtable}
}

As probabilidades somam \(0{,}3992+0{,}6008=1\). A probabilidade a priori de neve, \(P(Str=\text{snow\_covered})=0{,}10\), não multiplica a resposta, pois a rua coberta de neve já é uma condição conhecida. As tabelas de \(B\), \(K\) e \(Li\) também não entram: elas somam 1 ao serem marginalizadas, e não há evidência observada nesses nós.

\textbf{Resposta:} com as probabilidades adotadas, o dínamo apresenta tensão em uma rua coberta de neve com probabilidade de \textbf{39,92\%}. O resultado numérico depende das escolhas livres do item (c).


\bigskip
{\small
\textcolor{accent}{\textbf{Fonte do exercício}}\par
\textit{IA-Atividade\_Metodologias\_Ativas-RaciocinioProb.pdf}, questão 1. Universidade Federal do Amazonas, Instituto de Computação. Disciplina Inteligência Artificial, 2\textsuperscript{o} Trabalho, professor Edjard Mota, 11/10/2026.
}
\end{document}
